
\documentclass[aspectratio=169]{beamer}

\mode<presentation>
\usetheme{Boadilla}
\definecolor{NTHU1}{RGB}{127,16,132}
\definecolor{NTHU2}{RGB}{228,210,231}
\definecolor{blue}{RGB}{30,90,205}
\definecolor{red}{RGB}{213,94,0}
\definecolor{green}{RGB}{0,128,0}
\definecolor{charcoalgray}{RGB}{108,104,104}
\setbeamercolor{title}{fg=NTHU1}
\setbeamercolor{frametitle}{fg=NTHU1}
\setbeamercolor{block title}{bg=NTHU1, fg=white}
\setbeamercolor{block body}{bg=NTHU2}
\setbeamercolor{structure}{fg=NTHU1}
\setbeamercolor{item projected}{fg=white}
\setbeamercolor{item}{fg=NTHU1}
\setbeamercolor{subitem}{fg=NTHU1}
\setbeamercolor{section in toc}{fg=NTHU1}
\setbeamercolor{description item}{fg=NTHU1}
\setbeamercolor{caption name}{fg=NTHU1}
\setbeamercolor{button}{bg=NTHU1, fg=white}
\setbeamercolor{caption name}{fg=NTHU1}
\usepackage{graphics}
\usepackage{tikz}
\usepackage{amsmath}
\usepackage{bbm}
\usetikzlibrary{decorations.pathreplacing}
\usepackage{geometry}
\usepackage{booktabs}
\usepackage{bookmark}
\usepackage{multirow, makecell}
\usepackage{float}
\usepackage{fancyvrb}
\usepackage{caption}
\captionsetup[figure]{singlelinecheck = false, format= hang, justification=centering}
\captionsetup[subfigure]{singlelinecheck = false, format= hang, justification=raggedright}
\usepackage{subcaption}
\usepackage{adjustbox}
\usepackage{hyperref}
\usepackage{threeparttable}
\usepackage{appendixnumberbeamer} 
\usepackage[T1]{fontenc}
\usepackage[default]{lato} 
\usepackage{smartdiagram}
\smartdiagramset{
  back arrow disabled = true,
  module minimum width=1cm,
  module x sep = 3,
  uniform arrow color=true,
}
\usepackage{xcolor}
\usepackage{colortbl}
  \definecolor{light-gray}{gray}{0.99}

\newenvironment{wideitemize}{\itemize\addtolength{\itemsep}{10pt}}{\enditemize}
\newenvironment{wideenumerate}{\enumerate\addtolength{\itemsep}{10pt}}{\endenumerate}
\newenvironment{widedescription}{\description\addtolength{\itemsep}{10pt}}{\enddescription}
\hypersetup{
colorlinks=true,
linkcolor=NTHU1,
filecolor=green, 
urlcolor=blue,
}
\beamertemplatenavigationsymbolsempty
\setbeamercolor{author in head/foot}{bg=white, fg=NTHU1}
\setbeamercolor{title in head/foot}{bg=white, fg=NTHU1}
\setbeamercolor{date in head/foot}{bg=white, fg=NTHU1}
\setbeamercolor{section in head/foot}{bg=white, fg=NTHU1}
\setbeamercolor{headline}{bg=white}
\setbeamertemplate{footline}{
    \leavevmode%
    \hbox{%
        \begin{beamercolorbox}[wd=.333333\paperwidth,ht=2.25ex,dp=1ex,center]{date in head/foot}%
            \usebeamerfont{date in head/foot}%\insertdate
        \end{beamercolorbox}%
        \begin{beamercolorbox}[wd=.333333\paperwidth,ht=2.25ex,dp=1ex,center]{title in head/foot}%
            \usebeamerfont{title in head/foot}\insertshorttitle
        \end{beamercolorbox}%
        \begin{beamercolorbox}[wd=.333333\paperwidth,ht=2.25ex,dp=1ex,right]{date in head/foot}%
            \usebeamerfont{date in head/foot}\insertframenumber{}\hspace*{2em}
        \end{beamercolorbox}}%
        \vskip0pt%
    }
    %% May need to script the introduction!!!!!!
%\setbeamercolor{page number in head/foot}{fg=NTHU1}
\setbeamertemplate{section in toc}[sections numbered]
\setbeamertemplate{subsection in toc}{\leavevmode\leftskip=3em\rlap{\hskip-1.75em\inserttocsectionnumber.\inserttocsubsectionnumber}
\inserttocsubsection\par}
\setbeamerfont{subsection in toc}{size=\footnotesize}
%\setbeamertemplate{headline}{%
  %\begin{beamercolorbox}[ht=5.5ex]{section in head/foot}
    %\vskip2pt\insertnavigation{0.33\paperwidth}\vskip2pt
  %\end{beamercolorbox}%
%}
\newenvironment{transitionframe}{\setbeamercolor{background canvas}{bg=white}\setbeamertemplate{footline}{} \begin{frame}[noframenumbering]}{\end{frame}}

\makeatletter
\let\@@magyar@captionfix\relax
\makeatother

\title[]{Conflict and Peace} % Change this regularly
\subtitle[]{Week 10: Resistance - Overcoming collective action problems}
\author[Lee]{Seung-hun Lee }
\institute[]{Taipei School of Economics\\ National Tsing Hua University}

\date[]{November 7th, 2025}

\begin{document}
\begin{transitionframe}
\titlepage
\end{transitionframe}

%%% Color slides for section headers: Use for colloquium version (The ones bewteen \iffals and \fi)

\section{Overview of the topic}
\begin{frame}
\frametitle{What is collective action?} 
Definition of collective action
\begin{itemize}\setlength\itemsep{3pt}
\item Action taken by a group in pursuit of perceived shared interests
\begin{itemize}\setlength\itemsep{3pt}
\item Self-organizing protests to advocate certain causes and interests
\item Fighting in the military during wartimes
\item Negotiations / strikes carried out by labor unions (or any interest groups)
\item Management of communal resources
\end{itemize}
\item Involves sharing information and coordination among individuals to pool resources
\item Requires motivation and incentives from many individuals sharing common goals
\end{itemize}\medskip
This is a conceptual framework often used across many disciplines in social sciences
\begin{itemize}\setlength\itemsep{3pt}
\item Politics: How do interest groups sway votes? How do alliances sustain themselves?
\item Economics: How are public goods managed? How do groups organize?
\end{itemize}
\end{frame}

\begin{frame}
\frametitle{Why is it difficult to achieve?} 
Individual and collective incentives are misaligned
\begin{itemize}\setlength\itemsep{3pt}
\item Larger groups: Larger misalignment and smaller return per capita
\begin{itemize}
\item Coordination challenges: Larger coordination costs and information costs
\end{itemize}
\item Free-riding: Attempt to benefit from public good without individual contribution
\begin{itemize}\setlength\itemsep{3pt}
\item Game theory: If everyone thinks the same, public good is underprovided
\item Tragedy of commons: Rational individuals end up depleting shared resource
\end{itemize}
\end{itemize}\medskip
So how to overcome this problem?
\begin{itemize}\setlength\itemsep{3pt}
\item Decrease coordination costs: Communications and dispute settlement mechanisms
\item Align individual and collective incentives
\begin{itemize}\setlength\itemsep{3pt}
\item Make deviations costly: Monitoring, enforcement of contracts / laws
\item Internalize collective benefits: Increase awareness and salience of collective goals
\end{itemize}
\item Institutional setup to aggregate collective interests and costs effectively (governance)
\end{itemize}
\end{frame}

\begin{frame}
\frametitle{However, collective actions do happen worldwide} 
\begin{figure}
  \centering
  % First row
  \begin{subfigure}{0.4\textwidth}
    \includegraphics[width=0.75\linewidth, keepaspectratio]{fig1.jpg}
    \caption{Political protests}
  \end{subfigure}
  \begin{subfigure}{0.4\textwidth}
    \includegraphics[width=0.75\linewidth]{fig2.jpg}
    \caption{Environmental protests}
  \end{subfigure}
  
  % Second row
  \begin{subfigure}{0.4\textwidth}
    \includegraphics[width=0.75\linewidth]{fig3.jpg}
    \caption{Community service}
  \end{subfigure}
  \begin{subfigure}{0.4\textwidth}
    \includegraphics[width=0.75\linewidth]{fig4.jpg}
    \caption{Paying taxes}
  \end{subfigure}

\end{figure}
\end{frame}

\begin{frame}
\frametitle{Access to media and collective action problems} 
Academic interest in media increasing
\begin{itemize}\setlength\itemsep{3pt}
\item Dissemination of information and mechanism for improving accountability
\begin{itemize}\setlength\itemsep{3pt}
  \item Voters can monitor politicians (i.e. align interests of principals and agents)
  \item Improved transparency $\to$  make leaders behave well and foster trust in institutions
  \item Better information as basis for electoral outcomes
\end{itemize}
\item Social media, can help people mobilize
\begin{itemize}\setlength\itemsep{3pt}
  \item Inform users of the reality more accurately (information channel)
  \item Used to coordinate and mobilize group activities (coordination channel)
\end{itemize}
\end{itemize}\medskip
Just as these positives, media can work in reverse: Thus, an active arena of research
\begin{itemize}\setlength\itemsep{3pt}
  \item Potentially used as tools of repression (if government control is strong)
  \item Unchecked dissemination of misinformation
  \item Decaying trust in message may erode media's potential to lower coordination costs
\end{itemize}
\end{frame}

\begin{frame}
\frametitle{Access to media and collective action problems} 
\begin{figure}
\includegraphics[keepaspectratio,width=0.67\textwidth]{fig5.png}
\end{figure}
\end{frame}


\begin{transitionframe}
\frametitle{Roadmap for today}
\tableofcontents
\end{transitionframe}

\section{Understanding collective action - costs, leaders, and engagement}
\subsection{Ostrom (2000): Collective Action and the Evolution of Social Norms}
\begin{transitionframe}
\begin{center}
  \begin{huge}
    \textcolor{NTHU1}{%
Ostrom (2000): \\ Collective Action and the Evolution of Social Norms
    }
  \end{huge}
\end{center}
\end{transitionframe}

\begin{frame}
\frametitle{Q: What explains how collective action takes place}
Classic theories on when collective action happens 
\begin{itemize}\setlength\itemsep{3pt}
  \item Traditional thought: Groups take collective action whenever members jointly benefited
\item Mancur Olson (1965): Zero contribution thesis
\begin{quote}
  ``Unless the number of individuals in a group is quite \textcolor{blue}{small}, or unless there is \textcolor{blue}{coercion or some other special device} to make individuals act in their common interest, \textcolor{blue}{rational, self-interested} individuals will not act to achieve their common interests"
\end{quote}
\item In fairness, this explains some dilemma that we see in group action
\begin{itemize}\setlength\itemsep{3pt}
\item Why some people free-ride on collective benefits
\item Why small interest group yield more influence than general public
\end{itemize}
\item However, there are many incidences where zero contribution thesis is contradicted
\begin{itemize}
\item Individuals voluntarily organized to provide mutual protection and monitor each other
\end{itemize}
\item So what are the key factors that affect the likelihood of successful collective action?
\end{itemize}
\end{frame}

\begin{frame}
\frametitle{What are the existing lab evidence: Rational egoist perspective}
Behavior of self-interested rational individuals (Rational egoist) in public goods game
\begin{itemize}\setlength\itemsep{3pt}
  \item Individuals split their income into public good contribution and private consumption
  \item There are large number of individuals $N$ in this game
  \item Public goods yield some benefits for the individuals (i.e. individuals enjoy some share)
  \item But it reduces private consumption
\end{itemize}\medskip
Rational egoist predicts that equilibrium contribution is 0 (in most cases)
\begin{itemize}\setlength\itemsep{3pt}
  \item Given others' contribution, individuals access benefits without paying costs
  \item Since every individual thinks the same way, no one contributes
  \item Difficult to deviate from this equilibrium
  \begin{itemize}\setlength\itemsep{3pt}
  \item Suppose A contributes nonzero amount when everyone else contribute 0
  \item Unless the returns to individual from public good $>$ contribution, A is worse off
  \item And if the returns from public good $>$ contribution, equilibrium is not zero
\end{itemize}
\end{itemize}
\end{frame}

\begin{frame}
\frametitle{What are the existing lab evidence: What is actually found}
Findings from multiple lab experiments differ from rational egoist framework
\begin{itemize}\setlength\itemsep{3pt}
  \item Individuals contribute 40-60\% of endowments in finite games
  \item Contributions decay, but well above zero except for the final finite stage game
  \item Those who believe others will cooperate contribute to public goods themselves
  \item Individuals do learn to cooperate
  \item Face-to-Face communication increases cooperation substantially
  \item Individuals are willing to monitor/punish those who make below average Contributions
  \item Contextual framing matters (e.g. specifying sanctions, communications etc) 
\end{itemize} \medskip
These findings suggest that there may be more type of people than mere rational egoists
\end{frame}

\begin{frame}
\frametitle{Rebuilding collective action theory: Those who use norms}
Maybe there are some `conditional cooperators'
\begin{itemize}\setlength\itemsep{3pt}
  \item Those willing to cooperate when they estimate sufficient share of others also do
  \item Explains large contributions at initial stage (and can motivate rational egoists to join)
  \item Vary in tolerance for free-riding - the most tolerant ones contribute to the end
\end{itemize} \medskip
Or there are also `willing punishers'
\begin{itemize}\setlength\itemsep{3pt}
  \item Those willing to rebuke or impose costs on non-cooperators 
  \item Can reward active cooperators too
\end{itemize} \medskip
Emergence of these types?
\begin{itemize}\setlength\itemsep{3pt}
  \item Evolutionary process: Norms on collective action $>$ individual returns for survival 
  \item Thus, preferences adopt: People internalize cooperation as well as individual returns
\end{itemize} \medskip
\end{frame}

\begin{frame}
\frametitle{Common goods as collective action problem}
Common goods: How to cooperate when it is individually beneficial to use up resources
\begin{itemize}\setlength\itemsep{3pt}
  \item Nonexcludable: Difficult to prevent beneficiaries from enjoying benefits
  \item Rivalrous: Individual consumption subtracts from benefits available to others
  \item Think of a fish stock at international waters
  \item \textit{Note: Public goods are nonexcludable and \textbf{nonrivalrous}}
\end{itemize} \medskip
Management of these situations suggests rise of cooperation in collective action problems
\begin{itemize}\setlength\itemsep{3pt}
  \item Leader who oversees joint outcomes as key stimulus
  \item Clear boundary, well-tailored local rules

  \item Open to participation and communication in redistribution

\item Existence of accountable monitors and graduated sanctions help sustain collective action
\item Resolving conflicts through accessible low-cost, local arenas for discussion
\item Multilayer governance and minimal recognition of right to organize (not unanimity)
\end{itemize}
\end{frame}

\begin{frame}
\frametitle{Threats to sustained collective action}
What makes sustained collective action vulnerable?
\begin{itemize}\setlength\itemsep{3pt}
\item Migration changes composition of those who cooperate and understand social norms
\begin{itemize}
\item If in-migrants integrate well, these problems can be mitigated (two-way effort involved)
\end{itemize}
\item Lack of large-scale institutional and governance arrangements to settle disputes
\item Increasing opportunistic behavior through corruption
\item Too much reliance on external sources that are unaware of local contexts
\item Failure in horizontal/intergenerational transmission of norms through communication
\item National governments imposing single set of rules on all units of governance
\item Not adapting well to technological change
\end{itemize}\medskip
Growing need to understand how contextual factors affect cooperation and sustain collective action
\end{frame}

\subsection{Dippel, Heblich (2021): Leadership in Social Movements: Evidence from the ``Forty-Eighters" in Civil War}
\begin{transitionframe}
\begin{center}
  \begin{huge}
    \textcolor{NTHU1}{%
Dippel, Heblich (2021): \\ Leadership in Social Movements: Evidence from the ``Forty-Eighters" in Civil War
    }
  \end{huge}
\end{center}
\end{transitionframe}

\begin{frame}
\frametitle{Q: How do leaders contribute to social movements?} 
How do you enlist people to join a costly endeavor?
\begin{itemize}\setlength\itemsep{3pt}
\item Economic motives might be weak: Low and irregular pay
\item Legal monitoring may sometimes be absent (especially wartime)
\item Yet, we see many contemporary and historical examples of joining a costly endeavor
\item So what brings these people to this cause?
\end{itemize}\medskip
This paper uses introduction of Forty-Eighters to study effects of leaders on enlistments
\begin{itemize}\setlength\itemsep{3pt}
\item German revolutionaries who supported egalitarian causes expelled for their activities
\item Began campaigning for anti-slavery after 1854 (Kansas-Nebraska Act)
\item Huge campaigners against slavery and mobilizers of Union Army volunteers
\item Influence participation through local newspapers and social clubs
\end{itemize}\medskip
\end{frame}

\begin{frame}
\frametitle{How people fought in the US Civil War 1861-1865} 
US civil war
\begin{itemize}\setlength\itemsep{3pt}
 \item Fought over slavery and Southern secession 
\item 1/5 Adult men involved in fighting, 620,000 died in battle
\end{itemize}\medskip
So how did Germans get involved?
\begin{itemize}\setlength\itemsep{3pt}
\item German revolution (1848-1849): Many protested for republican and egalitarian causes
\item Systematically prosecuted and expelled to US - hence called `Forty-Eighters'
\item Many Forty-Eighters argued for emancipation following Kansas-Nebraska Act (1854)
\begin{itemize}
\item Slavery adopted based on popular sovereignty (raising salience of slavery)
\end{itemize}
\item Circulated arguments and swayed votes against slavery, later involved in enlistments
\begin{itemize}
\item Newspapers, Social clubs, public speeches, enlisted themselves to Union Army
\end{itemize}
\end{itemize}\medskip
\end{frame}

\begin{frame}
\frametitle{Data and empirical strategy} 
Data sources (town-level data)
\begin{itemize}\setlength\itemsep{3pt}
 \item Forty-Eighters: 493 Individual Biographies and genealogical source
 \begin{itemize}
\item Identify locations of their residence by 1856
  \end{itemize}
 \item Town characteristics from 1850, 1860 Census extracts
 \item Union Army enlistment: Digitalized Union Army Registers
\end{itemize}\medskip
Compare Forty-Eighters towns matched with similar non-Forty-Eighters towns
 \begin{itemize}\setlength\itemsep{3pt}
\item Match control group towns that share similar distribution of covariates 
 \begin{itemize}
  \item This is called propensity score matching
\end{itemize}
\item Compares the matched towns with this equation
\[
Y_i = \beta D[\text{FortyEighter}_i] + \gamma X_i +\phi_s+ \epsilon_i
\]\vspace{-20pt}
\item Also complemented with IV: Quasi-random destination distribution of Forty-Eighters
\begin{itemize}
\item Not enough time to select own destination, so followed co-passengers' destination
\end{itemize}
\end{itemize}
\end{frame} 

\begin{frame}
\frametitle{Enlistment and Forty-Eighters settlements match closely} 
\begin{figure}
\begin{subfigure}{0.495\textwidth}
\includegraphics[keepaspectratio, width=\textwidth]{dh2021_f1.png}
\caption{Forty-Eighters settlements}
\end{subfigure}
\begin{subfigure}{0.495\textwidth}
\includegraphics[keepaspectratio, width=\textwidth]{dh2021_f2.png}
\caption{Enlistments in Union Army}
\end{subfigure}

\end{figure}
\end{frame}

\begin{frame}
\frametitle{Enlistment and Forty-Eighters settlements match closely} 
\begin{figure}
\includegraphics[keepaspectratio, width=\textwidth]{dh2021_t1.png}
\end{figure}
\end{frame}

\begin{frame}
\frametitle{Newspaper circulation and social clubs explain the effects} 
\begin{figure}
\includegraphics[keepaspectratio, width=\textwidth]{dh2021_t2.png}
\end{figure}
\end{frame}

\begin{frame}
\frametitle{Takeaways and remaining questions} 
Other results
\begin{itemize}\setlength\itemsep{3pt}
\item IV results are consistent with previous results
\item Forty-eighters enlisted in Union Army themselves: Desertions $\Downarrow$ in those companies
\item Long-run: More likely to become members of groups advocating civil rights (NAACP)
\end{itemize}\medskip
Discussion
\begin{itemize}\setlength\itemsep{3pt}
\item Grassroots leadership can spread social movements
\item They serve as focal points in motivating belief through information
\item Emphasis on role of leaders are actively studied in Economics
\begin{itemize}\setlength\itemsep{3pt}
  \item How do politicians/leaders affect performances of their groups?
  \item Firm profits, economic growth, institutional stability etc
  \item Possible to use episodes of sudden absence to show how leaders matter

\end{itemize}
\end{itemize}
\end{frame}


\subsection{Bursztyin, Cantoni, Yang, Yuchtman, Zhang (2021): Persistent Political Engagement: Social Interactions and the Dynamics of Protest Movements [student presentation]}
\begin{transitionframe}
\begin{center}
  \begin{LARGE}
    \textcolor{NTHU1}{%
Bursztyin, Cantoni, Yang, Yuchtman, and  Zhang (2021): \\ Persistent Political Engagement: Social Interactions and the Dynamics of Protest Movements
 [student presentation]   }
  \end{LARGE}
\end{center}
\end{transitionframe}





\section{Role of media in resistance and collective action}
\subsection{Manacorda, Tesei (2020): Liberation Technology: Mobile Phones and Political Mobilization in Africa}
\begin{transitionframe}
\begin{center}
  \begin{huge}
    \textcolor{NTHU1}{%
Manacorda, Tesei (2020): \\ Liberation Technology: Mobile Phones and Political Mobilization in Africa
    }
  \end{huge}
\end{center}
\end{transitionframe}

\begin{frame}
\frametitle{Q: Can digital information foster political mobilization} 
Are mobile technologies `liberation technologies'?
\begin{itemize}\setlength\itemsep{3pt}
\item Can they assist marginalized groups in overcoming collective action problems?
\begin{itemize}
\item Low-cost, decentralized, multi-way communication can foster activism and mobilization
\item Especially when formal political participation is blocked or repressed
\end{itemize}
\item However, this may not always increase mobilization
\begin{itemize}\setlength\itemsep{3pt}
\item It may strengthen government surveillance and repression
\item May increase government accountability through information and greater transparency
\end{itemize}
\end{itemize}\medskip
This paper analyzes whether mobile networks foster political mobilization in Africa
\begin{itemize}\setlength\itemsep{3pt}
\item  Rapid mobile network expansion in 1998-2012 despite limited landline infrastructure
\item Rise of mobilizations motivated by grievances (foot riots, civil unrest like Arab spring)
\item Utilize georeferenced data on mobile coverage and uprisings
\end{itemize}
\end{frame}

\begin{frame}
\frametitle{Development of networks and protests in Africa} 
How mobile phones helped protests in Africa
\begin{itemize}\setlength\itemsep{3pt}
\item Arab Spring spread by video of self-immolation of Tunisian vendor
\begin{itemize}\setlength\itemsep{3pt}
\item The vendor protested authorities’ seizure of his goods after refusing to pay bride 
\item Circulated online and eventually to Al Jazeera
\item Protests erupted, fueled by discontent with low living standards, corruptions
\end{itemize}
\item This spilled over to Egypt, where mobile penetration was highest in Africa
\item Egyption President Mubarak resigned, despite efforts to shutdown internet
\end{itemize} \medskip
Other evidence suggest that mobile phones channeled grievances to protests
\begin{itemize}\setlength\itemsep{3pt}
\item Food riots in 2007-2012 across 14 countries spread with mobile messages and video recordings of police repression
\item Mobile coverage on average was 60\% higher in countries with food riots
\end{itemize}
\end{frame}

\begin{frame}
\frametitle{Development of networks and protests in Africa} 
\begin{figure}
  \begin{subfigure}{0.4\textwidth}
    \centering
    \includegraphics[keepaspectratio, width=\textwidth]{mt2020_f1.png}
    \caption{Mobile network coverage}
  \end{subfigure}
  \hfill
  \begin{subfigure}{0.55\textwidth}
    \centering
    \includegraphics[keepaspectratio, width=\textwidth]{mt2020_f2.png}
    \caption{Continentwide Protests and GDP}
  \end{subfigure} 
\end{figure}
\end{frame}


\begin{frame}
\frametitle{Data and empirical strategy} 
Data sources (Georeferenced unit: 55 km $\times$ 55 km grid)
\begin{itemize}\setlength\itemsep{3pt}
\item Mobile coverage: GSMA dataset of 2G, 3G, 4G coverage (1998-2012)
\item Protests $\to$ GDELT, ACLED, SCAD,  Participation $\to$ Afrobarometer surveys
\item Control for population, nightlights, rainfall, temperature, natural resources etc
\end{itemize}\medskip 
Leverage geographical (cell $j$ in country $c$) and yearly ($t$) variation in coverage and protests
\[
y_{jct} = \beta_1 \text{Cov}_{jct} + \beta_2\text{Cov}_{jct}\times\Delta \text{GDP}_{ct}+\gamma X_{jct} + f_j + f_{ct}+u_{jct}
\]\vspace{-20pt}
\begin{itemize}\setlength\itemsep{3pt}
\item However, mobile coverage may not be random - instrument $\text{Cov}_{jct}$ with lighting strikes
\begin{itemize}
\item Frequent lightning damages mobile networks(negatively correlated instrument)
\end{itemize}
\item Thus the first stages are
\[
\begin{aligned}
\text{Cov}_{jct} &= \delta_0+ \delta_1 \text{ltn}_{jc}\times t + \delta_2\text{ltn}_{jct}\times\Delta \text{GDP}_{ct}+\gamma_1 X_{jct} + f_j + f_{ct}+e_{jct}\\
\Delta\text{GDP}_{ct}\text{Cov}_{jct} &= \theta_0+ \theta_1 \text{ltn}_{jc}\times t + \theta_2\text{ltn}_{jct}\times\Delta \text{GDP}_{ct}+\gamma_2 X_{jct} + f_j + f_{ct}+\epsilon_{jct}
\end{aligned}
\]
\end{itemize}
\end{frame}

\begin{frame}
\frametitle{Protests in covered areas respond to \textit{economic shocks}} 
\begin{figure}
    \includegraphics[keepaspectratio, width=\textwidth]{mt2020_t1.png}
\end{figure}
\end{frame}

\begin{frame}
\frametitle{Heterogeneities across different contexts} 
\begin{figure}
    \includegraphics[keepaspectratio, width=0.7\textwidth]{mt2020_t2.png}
\end{figure}
\begin{itemize}
\item Mobilization $\Uparrow$ in autocracies, captured media, urban area, and areas with past conflicts
\end{itemize}
\end{frame}

\begin{frame}
\frametitle{Takeaways and discussion} 
Mobile phones help mobilization, but primarily through channeling grievances
\begin{itemize}\setlength\itemsep{3pt}
\item When grievances do not arise, effects are insignificant
\item Effects of mobile coverage \textit{interacted with GDP (economic downturns)} significant
\item Effects replicated at individual level participation outcomes
\item Thus, mobile phones facilitate mobilization by channeling grievances
\end{itemize}\medskip
Discussions and remaining questions
\begin{itemize}\setlength\itemsep{3pt}
\item  Mobile phones decrease cost of acquiring information and improves coordination
\begin{itemize}\setlength\itemsep{3pt}
\item Required for mobilization and other collective actions
\end{itemize}
\item Results generalizable beyond timeframe and region?
\begin{itemize}\setlength\itemsep{3pt}
\item What about states with advanced capacities and high technological level?
\end{itemize}
\end{itemize}\medskip
\end{frame}

\subsection{Guriev, Melnikov, Zhuravskaya (2021): 3G Internet and Confidence in Government}
\begin{transitionframe}
\begin{center}
  \begin{huge}
    \textcolor{NTHU1}{%
  Guriev, Melnikov, Zhuravskaya (2021): \\ 3G Internet and Confidence in Government
    }
  \end{huge}
\end{center}
\end{transitionframe}

\begin{frame}
\frametitle{Q: How does mobile network affect confidence in government?}
 Independent political thinking vs empowerment of nondemocratic regimes?
\begin{itemize}\setlength\itemsep{3pt}
\item Increase awareness and overcome collective action problems 
\item Easier dissemination of fake news, propaganda, and surveillance
\begin{itemize}
\item Former weakens confidence in government, latter strengthens it (esp nondemocracies)
\end{itemize}
\end{itemize}\medskip
This paper analyzes how expansion of 3G networks affect confidence in government
\begin{itemize}\setlength\itemsep{3pt}
\item 840,000 individuals across 2,232 regions in 116 countries from 2008-2017
\item 3G $\Uparrow \to$ Support for government $\Downarrow$ 
\begin{itemize}\setlength\itemsep{3pt}
\item Awareness of corruption $\Uparrow$, less trusting in institutions
\item Network channels outrage against status quo contagiously
\end{itemize}
\item Effect dominant when internet is not censored and traditional media is censored
\item Makes people aware of misgovernance: Effect present in highly corrupt countries
\end{itemize}
\end{frame}

\begin{frame}
\frametitle{Data and empirical strategy} 
Data sources
\begin{itemize}\setlength\itemsep{3pt}
\item Government approval: Gallup World Poll (2008-2017) 
\begin{itemize}
\item Yes/No on confidence in national government, judiciary, elections, corruption perception
\end{itemize}
\item 3G coverage: GSM association (same as Manacorda and Tesei (2020))
\item Institutional quality: Global Incidence of Corruption Index, Panama Papers Database (by journalists at ICIJ), Freedom House censorship index
\item Lightning: World Wide Lightning Location Network (WWLLN)  
\end{itemize}\medskip
Empirical strategy leverages 3G rollout variation across time and regions
\[ 
Y_{irt}=\gamma_1 \text{3G}_{rt} + \gamma_2 \text{Development}_{rt}+\lambda X_{irt} + \phi_r + \delta_t + \epsilon_{irt}
\]
\vspace{-20pt}
\begin{itemize}
  \item Lightning used to instrument 3G coverage
  \[
  \text{3G}_{rt} =  \beta_1 \text{ltn}_{r}\times t \times \text{Rich}_{c_r} + \beta_2\text{ltn}_{rt}\times \text{Poor}_{c_r} + \mu X_{irt} + \phi_r + \delta_t + \nu_{irt}
  \]
\end{itemize}
\end{frame}

\begin{frame}
\frametitle{Expansion of 3G networks, 2007-2018} 
\begin{figure}
  \centering
    \includegraphics[keepaspectratio, width=\textwidth]{gmz2021_f1.png}
    \caption{3G network expansion in 2007-2018}
\end{figure}
\end{frame}

\begin{frame}
\frametitle{Confidence in governance decreases across the board} 
\begin{figure}
  \centering
    \includegraphics[keepaspectratio, width=0.85\textwidth]{gmz2021_t1.png}
  
\end{figure}
\end{frame}

\begin{frame}
\frametitle{More so with uncensored internets and censored traditional media} 
\begin{figure}
  \centering
    \includegraphics[keepaspectratio, width=0.85\textwidth]{gmz2021_t2.png}
\end{figure}
\end{frame}

\begin{frame}
\frametitle{Decrease in approval driven by corruption detection} 
\begin{figure}
  \centering
    \includegraphics[keepaspectratio, width=0.9\textwidth]{gmz2021_t3.png}
\end{figure}
\end{frame}

\begin{frame}
\frametitle{Key takeaways and discussions} 
Major findings in this paper
\begin{itemize}\setlength\itemsep{3pt}
  \item 3G expansion decreases confidence in government
  \item Effects depends on censorship and corruption levels
\begin{itemize}\setlength\itemsep{3pt}
\item Internet censorship raises confidence, traditional media censorship decreases confidence
\item Effect stronger in more corrupt countries (by detecting new episodes of corruption)
\end{itemize}
\item Electoral consequences: Voters disillusioned with incumbents
\begin{itemize}
\item Vote share for incumbent $\downarrow$ Anti-establishment populists $\uparrow$
\end{itemize}
\end{itemize}\medskip
Discussions and remaining questions
\begin{itemize}\setlength\itemsep{3pt}
\item Who benefits from the 3G network?
\begin{itemize}\setlength\itemsep{3pt}
\item Nonpopulist oppositions did not gain vote shares
\item Anti-establishment includes extreme left and right-wing parties
\end{itemize}
\item What explains why anti-establishment are well off?
\begin{itemize}\setlength\itemsep{3pt}
\item What contents (potentially misinformation)?, Effects and mechanisms of misinformation?

\end{itemize}
\end{itemize}
\end{frame}

\subsection{Enikolopov, Marakin, Petrova (2020): Social Media and Protest Participation: Evidence from Russia}
\begin{transitionframe}
\begin{center}
  \begin{huge}
    \textcolor{NTHU1}{%
Enikolopov, Marakin, Petrova (2020): \\ Social Media and Protest Participation: Evidence from Russia
    }
  \end{huge}
\end{center}
\end{transitionframe}

\begin{frame}
\frametitle{Q: Does social media alleviate collective action problems?}
Overcoming collective action problems depends on ability to communicate 
\begin{itemize}\setlength\itemsep{3pt}
\item Social media allows conversations without intermediaries at low cost
\item This may lead to enhanced dissemination of information
\item However, if social media panders to rulers, information may decrease protests
\item Challenges: Selection bias in social media usage (individual, geographic characteristics)
\end{itemize}\medskip
This paper leverages unexpected protests and features of social media (VK) in Russia
\begin{itemize}\setlength\itemsep{3pt}
\item Leverages political protests in Russia in 2011
\item Protests more likely in areas where those with close ties to first users of VK lived
\item Distinguish information vs collective action channels
\begin{itemize}\setlength\itemsep{3pt}
\item Information: Increase awareness to those who do not know
\item Collective action: Facilitate coordination of logistics and peer pressure
\end{itemize}
\end{itemize}\medskip
\end{frame}

\begin{frame}
\frametitle{How social media facilitates protests and application to Russia} 
Two channels explaining effects of social media on protests
\begin{itemize}\setlength\itemsep{3pt}
\item Overall effect: Providing precise information + Assist in coordinating logistics
\begin{itemize}\setlength\itemsep{3pt}
\item Information: Negative content $\to$ Protest $\Uparrow$, Positive content $\to$ Protest $\Downarrow$
\item Coordination: Always raises protests
\end{itemize}
\item Coordination augmented by city size: Overcomes large coordination costs in cities 
\item If critical mass required: Protests occur only if certain threshold of users is reached
\end{itemize}\medskip
Russian protests in 2011 and Vkontakte (VK)
\begin{itemize}\setlength\itemsep{3pt}
\item Internet usage uncensored in 2011 (increased censoring after protests since 2012)
\item VK: First users personally approved by founder Pavel Durov - primarily SPSU students
\begin{itemize}
  \item Fun fact: Same person who found Telegram
\end{itemize}
\item Protests in 2011 against parliamentary election fraud in Dec 2011
\begin{itemize}
  \item VK users more aware of election fraud and coordinated protests
\end{itemize}
\end{itemize}
\end{frame}

\begin{frame}
\frametitle{It is the election where this happened...} 
\begin{figure}
\includegraphics[keepaspectratio,width=0.8\textwidth]{emp2020_f1.jpg}
\end{figure}
\end{frame}

\begin{frame}
\frametitle{Data sources and empirical strategy } 
Data sources (covering 625 cities in Russia)
\begin{itemize}\setlength\itemsep{3pt}
\item City of residence of VK users who joined before 2011 summer
\item Hand collected protest info (newspapers and online resources, 12/2012 - 05/2012)
\begin{itemize}\setlength\itemsep{3pt}
\item Complemented with alternative source from police and organizers 
\end{itemize}
\item VK founder cohort matched based on birth year, university, year of study (public)
\item Election data from Central Election Commission
\end{itemize}\medskip
Empirical strategy uses exogenous variation based on location of VK founder's cohort
\[
\begin{aligned}
  \text{OLS}: & Y_{i}= \beta_0 + \beta_1 VK_i + \beta_2 X_i + e_i\\
  \text{IV}: &VK_i=\gamma_0 + \gamma_1 f(\text{same SPSU cohort as VK founder})_i + \gamma_2 X_i+u_i\\
\end{aligned}
\]
\begin{itemize}\vspace{-10pt}
\item Verifies that same cohort as VK founder determines (not earlier or later) VK usage
\end{itemize}
\end{frame}

\begin{frame}
\frametitle{VK penetration, instrumented with VK founder cohort, raises protests} 
\begin{figure}
\includegraphics[keepaspectratio,width=0.8\textwidth]{emp2020_t1.png}
\end{figure}
\end{frame}

\begin{frame}
\frametitle{VK penetration effects larger in bigger cities} 
\begin{figure}
\includegraphics[keepaspectratio,width=\textwidth]{emp2020_f2.png}
\end{figure}
\end{frame}

\begin{frame}
\frametitle{Takeaways and remaining questions} 
Other notable results
\begin{itemize}\setlength\itemsep{3pt}
\item Electoral outcomes: VK penetration \textit{raised vote share} of ruling party
\begin{itemize}\setlength\itemsep{3pt}
\item On average, more positive contents describing ruling party documented
\item Voting outcomes only affected by information effect: You vote alone
\item Protests depends on information and coordination effects
\item Thus, in this case, information and coordination effects worked in opposite directions
\end{itemize}
\end{itemize}\medskip
Takeaways and remaining questions
\begin{itemize}\setlength\itemsep{3pt}
\item Social media does increase protests, primarily through coordination channel
\begin{itemize}\setlength\itemsep{3pt}
\item Information effect could still repress protests if government propaganda dominates
\end{itemize}
\item If different contents exist in social media, which one dominates?
\begin{itemize}\setlength\itemsep{3pt}
\item Could be a key to understanding misinformation propagation
\item Also necessary key to understand polarization on internet?
\end{itemize}

\end{itemize}
\end{frame}

\subsection{Gagliarducci, Onorato, Sobbrio, Tabellini (2020): War of the Waves: Radio and Resistance During World War II [student presentation]}
\begin{transitionframe}
\begin{center}
  \begin{huge}
    \textcolor{NTHU1}{%
Gagliarducci, Onorato, Sobbrio, Tabellini (2020): \\ War of the Waves: Radio and Resistance During World War II [student presentation]
    }
  \end{huge}
\end{center}
\end{transitionframe}

\begin{frame}
\frametitle{Takeaways and remaining questions} 
There are many ways to overcome collective action problems
\begin{itemize}\setlength\itemsep{3pt}
\item Role of leaders to stimulate a cause and organize a movements
\item Media's potential to bring individuals for a common cause
\begin{itemize}\setlength\itemsep{3pt}
\item Bring new and accurate information to those who were unaware of it previously
\item Lower the cost of coordination: Make it easy to reach out to people
\item Internet can achieve this if it is uncensored and traditional media is censored
\end{itemize}
\end{itemize}\medskip
Remaining questions
\begin{itemize}\setlength\itemsep{3pt}
\item Leadership can be used to motivate bad things (Tianyi Wang 2021 AER)
\begin{itemize}\setlength\itemsep{3pt}
\item Effective ways to hold decision-makers accountable is the key 
\end{itemize}
\item Likewise, media could channel negative influence
\begin{itemize}\setlength\itemsep{3pt}
\item Misinformation, echo-chamber effect (or both)
\item What are the ways to mitigate or control misinformation?
\end{itemize}
\end{itemize}
\end{frame}


%%%%%%%%%%%
\end{document}


